% ===========================================================================
% MECARAGE — Rapport de Projet de Fin d'Études
% Plateforme Intelligente de Gestion d'Atelier Mécanique
% ===========================================================================

\documentclass[12pt,a4paper]{report}

%% —— Encodage & Langue ——
\usepackage[utf8]{inputenc}
\usepackage[T1]{fontenc}
\usepackage[french]{babel}

%% —— Mise en page ——
\usepackage{geometry}
\geometry{a4paper,left=3cm,right=2.5cm,top=2.5cm,bottom=2.5cm,headheight=15pt}

%% —— Typographie ——
\usepackage{lmodern}
\usepackage{microtype}
\usepackage{setspace}
\onehalfspacing

%% —— Couleurs ——
\usepackage{xcolor}
\definecolor{indigo}{RGB}{79,70,229}
\definecolor{indigodark}{RGB}{55,48,163}
\definecolor{indigolite}{RGB}{199,210,254}
\definecolor{graylight}{RGB}{243,244,246}
\definecolor{prioH}{RGB}{220,38,38}
\definecolor{prioM}{RGB}{234,88,12}
\definecolor{prioL}{RGB}{22,163,74}
\definecolor{rowA}{RGB}{238,242,255}
\definecolor{rowB}{RGB}{255,255,255}

%% —— TikZ & PGF ——
\usepackage{tikz}
\usetikzlibrary{shapes.geometric,shapes.multipart,arrows.meta,positioning,
  fit,backgrounds,calc,shadows.blur,decorations.pathreplacing}
\usepackage{pgf-umlsd}
\usepackage{pgfplots}
\pgfplotsset{compat=1.18}

%% —— Graphiques & Flottants ——
\usepackage{graphicx}
\usepackage{float}
\usepackage{rotating}
\usepackage[labelfont=bf,labelsep=period,font=small]{caption}
\usepackage{subcaption}

%% —— Tableaux ——
\usepackage{booktabs}
\usepackage{array}
\usepackage{longtable}
\usepackage{multirow}
\usepackage{tabularx}
\usepackage{colortbl}

%% —— Listes & Code ——
\usepackage{enumitem}
\usepackage{listings}
\lstset{
  basicstyle=\footnotesize\ttfamily,
  breaklines=true,frame=tb,
  backgroundcolor=\color{graylight},
  keywordstyle=\color{indigo}\bfseries,
  commentstyle=\itshape\color{gray},
  numbers=left,numberstyle=\tiny\color{gray},numbersep=5pt}

%% —— En-têtes / Pieds de page ——
\usepackage{fancyhdr}
\pagestyle{fancy}
\fancyhf{}
\fancyhead[L]{\small\nouppercase{\leftmark}}
\fancyhead[R]{\small\textbf{MecaRage — PFE 2025/2026}}
\fancyfoot[C]{\thepage}
\renewcommand{\headrulewidth}{0.4pt}

%% —— Style des chapitres ——
\usepackage{titlesec}
\titleformat{\chapter}[display]
  {\normalfont\huge\bfseries\color{indigo}}
  {{\normalsize\bfseries\color{gray}Chapitre \thechapter}}{10pt}{\huge}
\titlespacing*{\chapter}{0pt}{-10pt}{25pt}
\titleformat{\section}{\large\bfseries\color{indigodark}}{\thesection}{1em}{}
\titleformat{\subsection}{\normalsize\bfseries\color{indigodark}}{\thesubsection}{1em}{}

%% —— Liens hypertextes ——
\usepackage{hyperref}
\hypersetup{
  colorlinks=true,linkcolor=indigo,citecolor=indigo,urlcolor=indigo,
  pdftitle={MecaRage -- Rapport PFE},pdfauthor={Equipe MecaRage},
  pdfkeywords={Angular,ASP.NET Core 9,MySQL,Firebase,Clean Architecture,CQRS}}

%% —— TikZ styles UML ——
\tikzset{
  umlclass/.style={
    rectangle split,rectangle split parts=3,
    rectangle split part fill={indigolite!60,white,white},
    draw=indigo,thick,font=\scriptsize,text width=4.2cm,align=left},
  ucboundary/.style={rectangle,draw=indigo,thick,rounded corners=5pt,
    fill=white!98,inner sep=12pt},
  ucactor/.style={rectangle,draw=indigo!60,fill=indigolite!30,rounded corners=3pt,
    font=\footnotesize\bfseries,text width=2.2cm,align=center,
    minimum height=0.8cm},
  ucell/.style={ellipse,draw=indigo,fill=indigolite!30,font=\scriptsize,
    text width=2.6cm,minimum height=0.75cm,align=center},
  ucline/.style={draw=gray!70,thin},
  classrel/.style={draw=indigo,thick,-{Stealth[length=6pt]}},
  classcomp/.style={draw=indigo,thick,{Diamond[length=6pt]}-{Stealth[length=6pt]}},
  classinh/.style={draw=indigo,thick,-{Triangle[open,length=8pt]}}}

\begin{document}

\pagenumbering{roman}
\setcounter{page}{1}

% ===========================================================================
% PAGE DE TITRE
% ===========================================================================
\begin{titlepage}
\centering
\vspace*{0.4cm}

{\large\scshape Université / École d'Ingénieur\par}
\vspace{0.2cm}
{\normalsize Département Informatique \& Génie Logiciel\par}
\vspace{0.3cm}
\noindent\rule{\linewidth}{1.5pt}
\vspace{0.3cm}

{\large Rapport de Projet de Fin d'Études (PFE)\par}
\vspace{0.8cm}

\begin{tikzpicture}
\node[rectangle,fill=indigo,rounded corners=8pt,
  text width=13cm,align=center,inner sep=1cm] {
  \color{white}
  \Huge\bfseries MecaRage\\[0.5cm]
  \Large Plateforme Intelligente de Gestion\\[0.2cm]
  d'Atelier Mécanique\\[0.4cm]
  \small Angular 21 \quad|\quad ASP.NET Core 9 \quad|\quad IA Gemini \quad|\quad Android Firebase
};
\end{tikzpicture}

\vspace{1.2cm}

\begin{minipage}{0.45\textwidth}
\begin{flushleft}
\textbf{Réalisé par~:}\\[0.3cm]
\hspace{0.5cm}[Prénom NOM]\\
\hspace{0.5cm}[Prénom NOM]
\end{flushleft}
\end{minipage}
\hfill
\begin{minipage}{0.45\textwidth}
\begin{flushright}
\textbf{Encadrant universitaire~:}\\[0.2cm]
\hspace{0.5cm}[Dr. Nom Encadrant]\\[0.4cm]
\textbf{Entreprise partenaire~:}\\[0.2cm]
\hspace{0.5cm}[Nom de l'entreprise]
\end{flushright}
\end{minipage}

\vfill
{\large Année Universitaire 2025--2026\par}
\end{titlepage}

% ===========================================================================
% RÉSUMÉ
% ===========================================================================
\chapter*{Résumé}
\addcontentsline{toc}{chapter}{Résumé}

\textbf{MecaRage} est une plateforme logicielle complète destinée à la gestion
moderne et intelligente des ateliers mécaniques. Elle repose sur une architecture
distribuée à trois couches~: un backend \textbf{ASP.NET Core 9} en Clean Architecture
(CQRS et MediatR), un frontend \textbf{Angular 21} avec Tailwind CSS, un service
d'intelligence artificielle fondé sur \textbf{FastAPI} et \textbf{Google Gemini}
pour le diagnostic contextuel par RAG (Retrieval-Augmented Generation), et un
module mobile \textbf{Android} (Java) intégrant \textbf{Firebase} pour les
notifications en temps réel.

La plateforme couvre l'intégralité du cycle de vie d'une intervention mécanique~:
dépôt des symptômes, diagnostic IA, prise de rendez-vous, affectation des mécaniciens,
rapport d'examen, facturation, approbation client, réparation, contrôle qualité et
paiement — le tout traçé par un suivi de cycle de vie unifié (\textit{Intervention
Lifecycle Tracker}).

\bigskip
\noindent\textbf{Mots-clés~:} Gestion d'atelier, Clean Architecture, CQRS,
Angular 21, ASP.NET Core 9, Firebase, Intelligence Artificielle, RAG, Android, MySQL.

\chapter*{Abstract}
\addcontentsline{toc}{chapter}{Abstract}

\textbf{MecaRage} is a comprehensive software platform for the modern and intelligent
management of mechanical workshops. It is built on a three-layer distributed
architecture~: an \textbf{ASP.NET Core 9} backend with Clean Architecture (CQRS and
MediatR), an \textbf{Angular 21} frontend with Tailwind CSS, an AI service based on
\textbf{FastAPI} and \textbf{Google Gemini} for RAG-based contextual diagnosis, and
an \textbf{Android} mobile module (Java) integrating \textbf{Firebase} for real-time
notifications.

The platform covers the complete lifecycle of a mechanical intervention~: symptom
submission, AI diagnosis, appointment scheduling, mechanic assignment, examination
report, invoicing, client approval, repair, quality validation and payment — all
tracked by a unified Intervention Lifecycle Tracker.

\bigskip
\noindent\textbf{Keywords~:} Workshop management, Clean Architecture, CQRS,
Angular 21, ASP.NET Core 9, Firebase, Artificial Intelligence, RAG, Android, MySQL.

% ===========================================================================
% TABLES
% ===========================================================================
\tableofcontents
\listoffigures
\listoftables

\cleardoublepage
\pagenumbering{arabic}
\setcounter{page}{1}

% ===========================================================================
% CHAPITRE 1 — LES BASES
% ===========================================================================
\chapter{Les Bases}

\section{Introduction}

Ce chapitre présente le contexte général du projet \textbf{MecaRage~:} la structure
de l'entreprise commanditaire, la problématique à l'origine du projet et la solution
proposée. Il s'agit d'établir un cadre de compréhension commun avant d'aborder les
aspects techniques et fonctionnels dans les chapitres suivants.

\section{Présentation de l'entreprise}

L'entreprise partenaire est spécialisée dans la gestion et l'exploitation de réseaux
d'ateliers mécaniques multi-sites. Elle coordonne plusieurs \textit{garages}
(centres de réparation) répartis géographiquement, chacun disposant d'une équipe de
mécaniciens supervisée par un chef d'atelier. La direction centrale (\textit{Admin
Entreprise}) assure la supervision de l'ensemble des opérations, des ressources
humaines et de la facturation.

Confrontée à une forte croissance de sa clientèle et à la multiplication de ses
ateliers, l'entreprise a identifié la nécessité de doter ses équipes d'un outil
numérique unifié pour remplacer ses processus manuels.

\subsection*{Organigramme des rôles métier}

\begin{center}
\begin{tikzpicture}[
  node distance=1.5cm,
  box/.style={rectangle,draw=indigo,fill=indigolite!30,rounded corners=3pt,
    font=\small\bfseries,text width=3.2cm,align=center,minimum height=0.8cm}]
  \node[box,fill=indigo!80,text=white] (sa) {Super Admin\\(Plateforme)};
  \node[box,below=of sa] (ae) {Admin Entreprise};
  \node[box,below left=1.2cm and 1.5cm of ae] (ca) {Chef d'Atelier};
  \node[box,below right=1.2cm and 1.5cm of ae] (cl) {Client};
  \node[box,below=1.2cm of ca] (mec) {Mécanicien};
  \draw[classrel] (sa) -- (ae);
  \draw[classrel] (ae) -- (ca);
  \draw[classrel] (ae) -- (cl);
  \draw[classrel] (ca) -- (mec);
\end{tikzpicture}
\end{center}

\section{Présentation du projet}

\subsection{Problématique}

La gestion quotidienne d'un réseau d'ateliers mécaniques implique une coordination
complexe entre plusieurs parties~: les clients qui déposent leurs véhicules, les
mécaniciens qui exécutent les travaux, le chef d'atelier qui orchestre les ressources,
et la direction qui supervise les finances. Cette coordination repose actuellement sur
des moyens peu adaptés~:

\begin{itemize}[noitemsep]
  \item Prise de rendez-vous par téléphone ou WhatsApp — pas de traçabilité.
  \item Rapports d'intervention écrits sur papier — risque de perte.
  \item Facturation produite sur tableur Excel — pas d'historique centralisé.
  \item Communication inter-équipes par messagerie informelle.
  \item Aucune visibilité en temps réel sur l'avancement des réparations.
  \item Aucun diagnostic préliminaire avant réception du véhicule.
\end{itemize}

\noindent Ces insuffisances entraînent des retards, des erreurs de facturation,
une mauvaise expérience client et une difficulté à piloter l'activité globalement.

\subsection{Étude de l'existant}

L'existant se compose d'une combinaison d'outils non intégrés~:

\begin{table}[H]
\centering
\caption{Outils existants et leurs limites}
\begin{tabularx}{\textwidth}{|l|X|X|}
\hline
\rowcolor{indigolite!30}
\textbf{Outil} & \textbf{Usage actuel} & \textbf{Limite principale} \\
\hline
WhatsApp / SMS & Prise de RDV, communications & Pas de traçabilité, informel \\
\hline
Excel & Inventaire pièces, facturation & Pas de collaboration temps réel \\
\hline
Papier & Rapports mécaniciens, bons & Perte, illisibilité, pas d'archivage \\
\hline
Tableau blanc & Suivi de l'atelier & Non numérique, local uniquement \\
\hline
\end{tabularx}
\end{table}

\subsection{Solution proposée}

\textbf{MecaRage} est une plateforme numérique unifiée qui digitalise l'ensemble du
cycle de vie d'une intervention mécanique. Elle offre~:

\begin{itemize}[noitemsep]
  \item Une \textbf{application web} (Angular 21) accessible par tous les rôles.
  \item Un \textbf{backend robuste} (ASP.NET Core 9) avec une architecture Clean
    Architecture et CQRS garantissant la maintenabilité.
  \item Un \textbf{service IA} (FastAPI + Google Gemini) produisant un diagnostic
    préliminaire par analyse sémantique des symptômes (RAG / BM25).
  \item Un \textbf{module mobile Android} (Java + Firebase) pour les clients et
    mécaniciens, avec notifications push en temps réel.
  \item Une \textbf{gestion multi-tenant} permettant à plusieurs entreprises
    d'exploiter la plateforme de façon isolée.
  \item Un \textbf{suivi de cycle de vie} automatique de chaque intervention, de
    l'examen jusqu'au paiement.
  \item Une \textbf{génération de factures PDF} via QuestPDF.
\end{itemize}

\section{Conclusion}

Ce chapitre a présenté le contexte métier qui a motivé la création de MecaRage.
Les lacunes identifiées dans les pratiques actuelles justifient pleinement la
conception d'une solution intégrée. Le chapitre suivant détaille les spécifications
fonctionnelles et techniques de la plateforme.

% ===========================================================================
% CHAPITRE 2 — SPÉCIFICATION FONCTIONNELLE ET TECHNIQUE
% ===========================================================================
\chapter{Spécification Fonctionnelle et Technique}

\section{Introduction}

Ce chapitre établit les exigences fonctionnelles et non-fonctionnelles de la
plateforme. Il présente le diagramme de cas d'utilisation global, le diagramme de
classes du domaine, le backlog produit complet et les choix techniques.

\section{Spécification Fonctionnelle}

\subsection{Exigences Fonctionnelles}

Les exigences fonctionnelles sont regroupées par acteur~:

\begin{longtable}{|p{0.5cm}|p{2.5cm}|p{8cm}|}
\caption{Exigences fonctionnelles par rôle}\label{tab:exigences}\\
\hline
\rowcolor{indigolite!60}
\textbf{N°} & \textbf{Acteur} & \textbf{Requirement} \\
\hline
\endfirsthead
\hline
\rowcolor{indigolite!60}
\textbf{N°} & \textbf{Acteur} & \textbf{Requirement} \\
\hline
\endhead
\hline
\endlastfoot
EF01 & Tous & S'authentifier avec email et mot de passe (JWT) \\
\hline
EF02 & Tous & Rafraîchir son token d'accès automatiquement \\
\hline
EF03 & Tous & Modifier son mot de passe \\
\hline
EF04 & SuperAdmin & Créer et gérer les tenants (entreprises) \\
\hline
EF05 & AdminEntreprise & Créer et gérer les garages du tenant \\
\hline
EF06 & AdminEntreprise & Créer et gérer le personnel (chefs, mécaniciens) \\
\hline
EF07 & AdminEntreprise & Consulter les statistiques (KPI) de l'atelier \\
\hline
EF08 & AdminEntreprise & Enregistrer le paiement d'une intervention \\
\hline
EF09 & Client & S'inscrire sur la plateforme \\
\hline
EF10 & Client & Gérer ses véhicules (ajout, consultation) \\
\hline
EF11 & Client & Soumettre un rapport de symptômes \\
\hline
EF12 & Système & Générer un diagnostic IA (Gemini RAG) automatiquement \\
\hline
EF13 & Client & Réserver un rendez-vous dans la plage proposée par le chef \\
\hline
EF14 & Client & Approuver ou refuser un devis (facture) \\
\hline
EF15 & Client & Suivre le cycle de vie de son intervention en temps réel \\
\hline
EF16 & Client & Télécharger sa facture en PDF \\
\hline
EF17 & ChefAtelier & Réviser les rapports de symptômes et proposer une plage RDV \\
\hline
EF18 & ChefAtelier & Approuver ou refuser des rendez-vous \\
\hline
EF19 & ChefAtelier & Assigner des mécaniciens à une tâche de réparation \\
\hline
EF20 & ChefAtelier & Réviser le rapport d'examen du mécanicien \\
\hline
EF21 & ChefAtelier & Créer et finaliser une facture (devis) \\
\hline
EF22 & ChefAtelier & Valider la réparation et notifier le client \\
\hline
EF23 & Mécanicien & Consulter ses tâches assignées \\
\hline
EF24 & Mécanicien & Soumettre un rapport d'examen (pièces, observations) \\
\hline
EF25 & Mécanicien & Soumettre un rapport de fin de réparation \\
\hline
EF26 & Système & Envoyer des notifications in-app (mysql) et push (FCM) \\
\hline
EF27 & AdminEntreprise & Gérer l'inventaire des pièces détachées \\
\hline
EF28 & Tous & Accéder via application Android (client + mécanicien) \\
\hline
\end{longtable}

\subsection{Diagramme de cas d'utilisation global}

\begin{figure}[H]
\centering
\resizebox{\textwidth}{!}{%
\begin{tikzpicture}[auto]
  % System boundary
  \node[ucboundary,minimum width=14cm,minimum height=18cm,
    label={[font=\large\bfseries\color{indigo}]above:Système MecaRage}]
    (sys) at (6,0) {};

  % ----- ACTORS -----
  \node[ucactor] (cli) at (-3.5, 2)  {Client};
  \node[ucactor] (mec) at (-3.5,-4)  {Mécanicien};
  \node[ucactor] (che) at (15.5, 4)  {Chef\\d'Atelier};
  \node[ucactor] (adm) at (15.5,-1)  {Admin\\Entreprise};
  \node[ucactor] (sad) at (15.5,-6)  {Super\\Admin};
  \node[ucactor,fill=gray!15] (ia)  at (6, 11)  {Service IA\\(Gemini)};

  % ----- USE CASES — Authentication -----
  \node[ucell] (uc01) at (3.5, 9.0) {Authentification};
  \node[ucell] (uc02) at (8.5, 9.0) {Gestion du\\Profil};

  % ----- USE CASES — Client -----
  \node[ucell] (uc03) at (2.0, 6.5)  {Soumettre Rapport\\de Symptômes};
  \node[ucell] (uc04) at (6.0, 6.5)  {Réserver\\Rendez-vous};
  \node[ucell] (uc05) at (10.0,6.5)  {Approuver /\\Refuser Devis};
  \node[ucell] (uc06) at (4.0, 4.5)  {Suivre\\Intervention};
  \node[ucell] (uc07) at (8.0, 4.5)  {Télécharger\\Facture PDF};

  % ----- USE CASES — Mechanic -----
  \node[ucell] (uc08) at (2.5, 2.0)  {Consulter Tâches\\Assignées};
  \node[ucell] (uc09) at (7.0, 2.0)  {Soumettre Rapport\\d'Examen};
  \node[ucell] (uc10) at (2.5,-0.5)  {Soumettre Rapport\\de Réparation};

  % ----- USE CASES — Chef -----
  \node[ucell] (uc11) at (10.0,2.0)  {Réviser Rapport\\Symptômes};
  \node[ucell] (uc12) at (10.0,-0.5) {Gérer\\Rendez-vous};
  \node[ucell] (uc13) at (7.0,-2.5)  {Assigner\\Mécaniciens};
  \node[ucell] (uc14) at (10.0,-4.0) {Créer et Finaliser\\Facture};
  \node[ucell] (uc15) at (7.0,-5.5)  {Valider Réparation\\et Notifier};

  % ----- USE CASES — Admin -----
  \node[ucell] (uc16) at (2.5,-4.0)  {Gérer Garages\\et Personnel};
  \node[ucell] (uc17) at (2.5,-6.5)  {Gérer Pièces\\Détachées};
  \node[ucell] (uc18) at (6.0,-7.5)  {Enregistrer\\Paiement};

  % ----- USE CASES — SuperAdmin -----
  \node[ucell] (uc19) at (10.0,-7.5) {Gérer Tenants};

  % ----- AI System -----
  \node[ucell] (uc20) at (6.0, 8.0)  {Diagnostic IA\\Automatique};

  % ----- ASSOCIATIONS — Client -----
  \draw[ucline] (cli.east) -- (uc01.west);
  \draw[ucline] (cli.east) -- (uc03.west);
  \draw[ucline] (cli.east) -- (uc04.west);
  \draw[ucline] (cli.east) -- (uc05.west);
  \draw[ucline] (cli.east) -- (uc06.west);
  \draw[ucline] (cli.east) -- (uc07.west);

  % ----- ASSOCIATIONS — Mechanic -----
  \draw[ucline] (mec.east) -- (uc01.west);
  \draw[ucline] (mec.east) -- (uc08.west);
  \draw[ucline] (mec.east) -- (uc09.west);
  \draw[ucline] (mec.east) -- (uc10.west);

  % ----- ASSOCIATIONS — Chef -----
  \draw[ucline] (che.west) -- (uc01.east);
  \draw[ucline] (che.west) -- (uc11.east);
  \draw[ucline] (che.west) -- (uc12.east);
  \draw[ucline] (che.west) -- (uc13.east);
  \draw[ucline] (che.west) -- (uc14.east);
  \draw[ucline] (che.west) -- (uc15.east);

  % ----- ASSOCIATIONS — Admin -----
  \draw[ucline] (adm.west) -- (uc01.east);
  \draw[ucline] (adm.west) -- (uc16.east);
  \draw[ucline] (adm.west) -- (uc17.east);
  \draw[ucline] (adm.west) -- (uc18.east);
  \draw[ucline] (adm.west) -- (uc14.east);

  % ----- ASSOCIATIONS — SuperAdmin -----
  \draw[ucline] (sad.west) -- (uc01.east);
  \draw[ucline] (sad.west) -- (uc19.east);

  % ----- ASSOCIATIONS — IA -----
  \draw[ucline,dashed] (ia.south) -- (uc20.north);
  \draw[ucline,dashed,-{Stealth}] (uc03.north) -- node[right,font=\scriptsize]
    {\texttt{<<include>>}} (uc20.south west);

\end{tikzpicture}
}%
\caption{Diagramme de cas d'utilisation global — MecaRage}
\label{fig:uc-global}
\end{figure}

\subsection{Diagramme de classes global}

\begin{sidewaysfigure}[p]
\centering
\resizebox{0.97\textheight}{!}{%
\begin{tikzpicture}[node distance=0.6cm and 1.5cm]

% Tenant
\node[umlclass] (tenant) {
  \textbf{Tenant}
  \nodepart{two}
  + id : Guid\\+ name : String\\+ slug : String\\+ email : String\\+ isActive : bool
  \nodepart{three}
  \textit{(aucune méthode publique)}
};

% Garage
\node[umlclass,right=3.5cm of tenant] (garage) {
  \textbf{Garage}
  \nodepart{two}
  + id : Guid\\+ tenantId : Guid\\+ name : String\\+ address : String\\+ city : String\\+ phone : String
  \nodepart{three}
  \textit{(aucune méthode publique)}
};

% User
\node[umlclass,below=2.5cm of tenant] (user) {
  \textbf{User}
  \nodepart{two}
  + id : Guid\\+ firstName : String\\+ email : String\\+ role : UserRole\\+ tenantId : Guid\\+ garageId : Guid?
  \nodepart{three}
  \textit{(aucune méthode publique)}
};

% Vehicle
\node[umlclass,right=3.5cm of user] (vehicle) {
  \textbf{Vehicle}
  \nodepart{two}
  + id : Guid\\+ clientId : Guid\\+ brand : String\\+ model : String\\+ year : int\\+ licensePlate : String
  \nodepart{three}
  \textit{(aucune méthode publique)}
};

% SymptomReport
\node[umlclass,right=3.5cm of vehicle] (symptom) {
  \textbf{SymptomReport}
  \nodepart{two}
  + id : Guid\\+ clientId : Guid\\+ vehicleId : Guid\\+ symptomsDesc : String\\+ aiPredictedIssue : String\\+ status : SymptomReportStatus
  \nodepart{three}
  \textit{(aucune méthode publique)}
};

% Appointment
\node[umlclass,below=2.5cm of vehicle] (appt) {
  \textbf{Appointment}
  \nodepart{two}
  + id : Guid\\+ clientId : Guid\\+ vehicleId : Guid\\+ garageId : Guid\\+ preferredDate : DateTime\\+ status : AppointmentStatus
  \nodepart{three}
  \textit{(aucune méthode publique)}
};

% SparePart
\node[umlclass,left=3.5cm of appt] (spare) {
  \textbf{SparePart}
  \nodepart{two}
  + id : Guid\\+ garageId : Guid\\+ code : String\\+ name : String\\+ unitPrice : decimal\\+ stockQty : int
  \nodepart{three}
  \textit{(aucune méthode publique)}
};

% Invoice
\node[umlclass,right=3.5cm of appt] (invoice) {
  \textbf{Invoice}
  \nodepart{two}
  + id : Guid\\+ appointmentId : Guid\\+ invoiceNumber : String\\+ serviceFee : decimal\\+ totalAmount : decimal\\+ status : InvoiceStatus
  \nodepart{three}
  \textit{(aucune méthode publique)}
};

% RepairTask
\node[umlclass,below=2.5cm of appt] (repair) {
  \textbf{RepairTask}
  \nodepart{two}
  + id : Guid\\+ appointmentId : Guid\\+ taskTitle : String\\+ status : RepairTaskStatus\\+ assignedAt : DateTime\\+ completionNotes : String?
  \nodepart{three}
  \textit{(aucune méthode publique)}
};

% Intervention
\node[umlclass,right=3.5cm of repair] (interv) {
  \textbf{Intervention}
  \nodepart{two}
  + id : Guid\\+ appointmentId : Guid?\\+ invoiceId : Guid?\\+ status : InterventionLifecycleStatus\\+ paymentAmount : decimal?\\+ paymentMethod : String?
  \nodepart{three}
  \textit{(aucune méthode publique)}
};

% Notification
\node[umlclass,left=3.5cm of repair] (notif) {
  \textbf{Notification}
  \nodepart{two}
  + id : Guid\\+ recipientId : Guid\\+ title : String\\+ message : String\\+ isRead : bool\\+ createdAt : DateTime
  \nodepart{three}
  \textit{(aucune méthode publique)}
};

% ----- RELATIONSHIPS -----
\draw[classrel] (garage.west) -- node[above,font=\scriptsize]{N}
  node[below,font=\scriptsize]{1} (tenant.east);
\draw[classrel] (user.north) -- node[right,font=\scriptsize]{N}
  node[left,font=\scriptsize]{1} (tenant.south);
\draw[classrel] (vehicle.west) -- node[above,font=\scriptsize]{N}
  node[below,font=\scriptsize]{1} (user.east);
\draw[classrel] (symptom.south) -- node[right,font=\scriptsize]{N}
  node[left,font=\scriptsize]{1} (vehicle.north);
\draw[classrel] (appt.north) -- node[right,font=\scriptsize]{N}
  node[left,font=\scriptsize]{0..1} (vehicle.south);
\draw[classrel] (appt.east) -- node[above,font=\scriptsize]{0..1}
  node[below,font=\scriptsize]{1} (invoice.west);
\draw[classrel] (appt.south) -- node[right,font=\scriptsize]{0..1}
  node[left,font=\scriptsize]{1} (repair.north);
\draw[classcomp] (repair.south) -- (interv.north west);
\draw[classcomp] (invoice.south) -- (interv.north east);
\draw[classrel] (spare.east) -- node[above,font=\scriptsize]{N}
  node[below,font=\scriptsize]{1} (garage.south west);
\draw[classrel] (notif.north east) -- node[above,font=\scriptsize]{N}
  node[below,font=\scriptsize]{1} (user.south west);

\end{tikzpicture}
}%
\caption{Diagramme de classes global — principales entités du domaine}
\label{fig:class-global}
\end{sidewaysfigure}

\subsection{Backlog Produit}

\begin{longtable}{|p{0.7cm}|p{4cm}|p{5.5cm}|p{1.3cm}|p{0.7cm}|p{1.2cm}|}
\caption{Backlog Produit — MecaRage}\label{tab:backlog}\\
\hline
\rowcolor{indigolite!60}
\textbf{ID} & \textbf{En tant que} & \textbf{Je veux} & \textbf{Priorité} & \textbf{SP} & \textbf{Sprint} \\
\hline
\endfirsthead
\hline
\rowcolor{indigolite!60}
\textbf{ID} & \textbf{En tant que} & \textbf{Je veux} & \textbf{Priorité} & \textbf{SP} & \textbf{Sprint} \\
\hline
\endhead
\hline
\endlastfoot
US01 & Utilisateur & M'inscrire avec email et mot de passe & \textcolor{prioH}{\bfseries Haute} & 3 & 1 \\
\hline
US02 & Utilisateur & Me connecter et recevoir un JWT & \textcolor{prioH}{\bfseries Haute} & 3 & 1 \\
\hline
US03 & Utilisateur & Rafraîchir mon token automatiquement & \textcolor{prioH}{\bfseries Haute} & 2 & 1 \\
\hline
US04 & Utilisateur & Modifier mon mot de passe & \textcolor{prioM}{\bfseries Moyenne} & 2 & 1 \\
\hline
US05 & Super Admin & Créer et gérer les tenants & \textcolor{prioH}{\bfseries Haute} & 3 & 2 \\
\hline
US06 & Admin Entreprise & Créer et gérer des garages & \textcolor{prioH}{\bfseries Haute} & 5 & 2 \\
\hline
US07 & Admin Entreprise & Gérer le personnel (chefs, mécaniciens) & \textcolor{prioH}{\bfseries Haute} & 5 & 2 \\
\hline
US08 & Client & Enregistrer et gérer mes véhicules & \textcolor{prioH}{\bfseries Haute} & 3 & 2 \\
\hline
US09 & Admin Entreprise & Gérer l'inventaire des pièces détachées & \textcolor{prioM}{\bfseries Moyenne} & 5 & 2 \\
\hline
US10 & Client & Soumettre un rapport de symptômes & \textcolor{prioH}{\bfseries Haute} & 5 & 3 \\
\hline
US11 & Système & Générer un diagnostic IA (RAG/Gemini) & \textcolor{prioH}{\bfseries Haute} & 8 & 3 \\
\hline
US12 & Chef d'Atelier & Réviser les rapports et proposer une plage RDV & \textcolor{prioH}{\bfseries Haute} & 5 & 3 \\
\hline
US13 & Client & Réserver un rendez-vous & \textcolor{prioH}{\bfseries Haute} & 3 & 3 \\
\hline
US14 & Chef d'Atelier & Approuver ou refuser un rendez-vous & \textcolor{prioH}{\bfseries Haute} & 3 & 3 \\
\hline
US15 & Chef d'Atelier & Assigner des mécaniciens à une tâche & \textcolor{prioH}{\bfseries Haute} & 5 & 4 \\
\hline
US16 & Mécanicien & Consulter mes tâches assignées & \textcolor{prioH}{\bfseries Haute} & 3 & 4 \\
\hline
US17 & Mécanicien & Soumettre un rapport d'examen & \textcolor{prioH}{\bfseries Haute} & 5 & 4 \\
\hline
US18 & Chef d'Atelier & Réviser et valider le rapport d'examen & \textcolor{prioH}{\bfseries Haute} & 3 & 4 \\
\hline
US19 & Chef d'Atelier & Créer et finaliser une facture (devis) & \textcolor{prioH}{\bfseries Haute} & 5 & 4 \\
\hline
US20 & Client & Approuver ou refuser un devis & \textcolor{prioH}{\bfseries Haute} & 3 & 4 \\
\hline
US21 & Mécanicien & Soumettre la fin de réparation & \textcolor{prioH}{\bfseries Haute} & 3 & 4 \\
\hline
US22 & Chef d'Atelier & Valider et notifier le client (véhicule prêt) & \textcolor{prioH}{\bfseries Haute} & 3 & 4 \\
\hline
US23 & Admin Entreprise & Enregistrer le paiement au retrait & \textcolor{prioH}{\bfseries Haute} & 3 & 4 \\
\hline
US24 & Système & Tracer le cycle de vie (Intervention Tracker) & \textcolor{prioH}{\bfseries Haute} & 8 & 4 \\
\hline
US25 & Tout utilisateur & Recevoir des notifications in-app et push & \textcolor{prioM}{\bfseries Moyenne} & 5 & 4 \\
\hline
US26 & Client/Mécanicien & Télécharger la facture en PDF & \textcolor{prioM}{\bfseries Moyenne} & 3 & 4 \\
\hline
US27 & Client/Mécanicien & Utiliser l'application Android native & \textcolor{prioM}{\bfseries Moyenne} & 13 & 4 \\
\hline
\end{longtable}

\section{Spécification Technique}

\subsection{Exigences Non-Fonctionnelles}

\begin{table}[H]
\centering
\caption{Exigences non-fonctionnelles}
\begin{tabularx}{\textwidth}{|l|X|}
\hline
\rowcolor{indigolite!30}
\textbf{Critère} & \textbf{Exigence} \\
\hline
Performance & Temps de réponse API $<$ 300 ms pour 95\% des requêtes \\
\hline
Sécurité & JWT HS256, hachage bcrypt, isolation multi-tenant, CORS contrôlé \\
\hline
Scalabilité & Architecture stateless; frontend découplé; services indépendamment déployables \\
\hline
Disponibilité & Déploiement via Docker Compose avec healthchecks et \texttt{restart: unless-stopped} \\
\hline
Maintenabilité & Clean Architecture, CQRS, tests unitaires xUnit, couverture $\geq 70\%$ \\
\hline
Compatibilité & Android 8+ (API 26); navigateurs modernes (Chrome, Firefox, Edge) \\
\hline
\end{tabularx}
\end{table}

\subsection{Outils et Environnement de Développement}

\begin{table}[H]
\centering
\caption{Stack technologique}
\begin{tabularx}{\textwidth}{|l|l|X|}
\hline
\rowcolor{indigolite!30}
\textbf{Couche} & \textbf{Technologie} & \textbf{Version / Rôle} \\
\hline
Frontend & Angular + Tailwind CSS & v21.1 / SPA réactive, dark mode \\
\hline
Backend & ASP.NET Core + MediatR & v9.0 / API REST, CQRS \\
\hline
Base de données & MySQL / Pomelo EF Core & 8.0 / ORM relationnel \\
\hline
Cache & Redis & 7.0 / sessions, invalidation \\
\hline
IA & Python FastAPI + Gemini & 3.13 / RAG BM25 + LLM \\
\hline
Mobile & Android Java + Firebase & SDK 26+ / RTDB + FCM \\
\hline
PDF & QuestPDF & 2026.2.4 / génération factures \\
\hline
Déploiement & Docker + Docker Compose & 25+ / containerisation \\
\hline
CI/CD & Ansible + playbooks & — / déploiement automatisé \\
\hline
Monitoring & Prometheus + Grafana & — / métriques et alertes \\
\hline
IDE & JetBrains Rider / VS Code & — / développement \\
\hline
\end{tabularx}
\end{table}

\section{Conclusion}

Ce chapitre a formalisé les besoins de la plateforme par 28 user stories organisées
en 4 sprints, et a présenté les diagrammes de cas d'utilisation et de classes du
domaine. Les chapitres suivants détaillent la conception et le développement de
chaque sprint.

% ===========================================================================
% CHAPITRE 3 — SPRINT 1 : AUTHENTIFICATION ET GESTION DES PROFILS
% ===========================================================================
\chapter{Sprint 1 : Authentification et Gestion des Profils}

\section{Introduction}

Ce premier sprint pose les fondations de la plateforme en implémentant le système
d'authentification et de gestion des profils. Il couvre l'inscription, la connexion,
la gestion des tokens JWT et le changement de mot de passe. Ce sprint constitue
le prérequis de toutes les fonctionnalités ultérieures.

\section{Identification du Backlog Sprint 1}

\begin{table}[H]
\centering
\caption{Backlog Sprint 1 — Authentification}
\begin{tabularx}{\textwidth}{|l|X|l|c|}
\hline
\rowcolor{indigolite!60}
\textbf{ID} & \textbf{User Story} & \textbf{Priorité} & \textbf{Story Points} \\
\hline
US01 & En tant qu'utilisateur, je veux m'inscrire avec mon email et mon mot de passe &
  \textcolor{prioH}{\bfseries Haute} & 3 \\
\hline
US02 & En tant qu'utilisateur, je veux me connecter et recevoir un token JWT &
  \textcolor{prioH}{\bfseries Haute} & 3 \\
\hline
US03 & En tant qu'utilisateur, je veux que mon token se rafraîchisse automatiquement &
  \textcolor{prioH}{\bfseries Haute} & 2 \\
\hline
US04 & En tant qu'utilisateur, je veux pouvoir modifier mon mot de passe &
  \textcolor{prioM}{\bfseries Moyenne} & 2 \\
\hline
\multicolumn{3}{|r|}{\textbf{Total Sprint 1}} & \textbf{10 SP} \\
\hline
\end{tabularx}
\end{table}

\noindent\textbf{Tâches techniques du Sprint 1~:}
\begin{enumerate}[noitemsep]
  \item Conception du modèle de domaine \texttt{User} et \texttt{Tenant}.
  \item Implémentation des commandes CQRS~: \texttt{RegisterCommand}, \texttt{LoginCommand},
    \texttt{RefreshTokenCommand}, \texttt{ChangePasswordCommand}.
  \item Mise en place du middleware JWT et du \texttt{ClaimsPrincipalExtensions}.
  \item Création des endpoints REST dans \texttt{AuthController}.
  \item Configuration de l'isolation multi-tenant par filtre global EF Core.
  \item Développement des écrans Angular~: \texttt{LoginComponent}, \texttt{SignupComponent}.
  \item Implémentation de l'intercepteur HTTP pour l'injection du Bearer Token.
\end{enumerate}

\section{Conception}

\subsection{Diagramme de classes — Sprint 1}

\begin{figure}[H]
\centering
\begin{tikzpicture}[node distance=0.8cm and 4cm]

\node[umlclass] (user) {
  \textbf{User}
  \nodepart{two}
  + id : Guid\\+ firstName : String\\+ lastName : String\\+ email : String\\+ passwordHash : String\\+ role : UserRole\\+ tenantId : Guid?\\+ garageId : Guid?\\+ refreshToken : String?\\+ refreshTokenExpiry : DateTime?\\+ isActive : bool
  \nodepart{three}
  (géré via commands)
};

\node[umlclass,right=4cm of user] (tenant) {
  \textbf{Tenant}
  \nodepart{two}
  + id : Guid\\+ name : String\\+ slug : String\\+ email : String\\+ phone : String\\+ isActive : bool
  \nodepart{three}
  (géré via commands)
};

\node[umlclass,below=2.5cm of user] (jwt) {
  \textbf{LoginResponse}
  \nodepart{two}
  + success : bool\\+ accessToken : String\\+ refreshToken : String\\+ expiresAt : DateTime\\+ userId : Guid\\+ role : String
  \nodepart{three}
  (DTO de retour)
};

\draw[classrel] (user.east) -- node[above,font=\scriptsize]{N}
  node[below,font=\scriptsize]{1} (tenant.west)
  node[pos=0.98,right,font=\scriptsize]{appartient à};

\draw[classrel,dashed] (user.south) -- node[right,font=\scriptsize]
  {\texttt{<<produce>>}} (jwt.north);

\end{tikzpicture}
\caption{Diagramme de classes — Sprint 1 (Authentification)}
\label{fig:class-sprint1}
\end{figure}

\noindent\textbf{Enumération \texttt{UserRole}~:}
\begin{itemize}[noitemsep]
  \item \texttt{SuperAdmin} — Administration globale de la plateforme
  \item \texttt{AdminEntreprise} — Gestion d'un tenant / entreprise
  \item \texttt{ChefAtelier} — Supervisor d'un garage
  \item \texttt{Mecanicien} — Technicien de réparation
  \item \texttt{Client} — Propriétaire de véhicule
\end{itemize}

\subsection{Diagrammes de séquence}

\subsubsection*{SC1 — Connexion (Login)}

\begin{figure}[H]
\centering
\resizebox{\textwidth}{!}{%
\begin{sequencediagram}
  \newthread{cli}{:Client}
  \newinst[1.5]{ctrl}{:AuthController}
  \newinst[1.5]{hand}{:LoginCmd\\Handler}
  \newinst[1.5]{jwt}{:JwtService}
  \newinst[1.5]{db}{:Database}

  \begin{call}{cli}{POST /api/auth/login\{email,password\}}{ctrl}{}
    \begin{call}{ctrl}{Handle(LoginCommand)}{hand}{}
      \begin{call}{hand}{FindByEmail(email)}{db}{User | null}
      \end{call}
      \begin{callself}{hand}{VerifyHash(password, hash)}{}
      \end{callself}
      \begin{call}{hand}{GenerateToken(user)}{jwt}{accessToken}
      \end{call}
      \begin{call}{hand}{GenerateRefreshToken()}{jwt}{refreshToken}
      \end{call}
      \begin{call}{hand}{SaveRefreshToken(userId, token)}{db}{}
      \end{call}
    \end{call}
    \mess{ctrl}{200 OK \{accessToken, refreshToken, role\}}{cli}
  \end{call}

  \begin{callself}{cli}{StoreToken(localStorage)}{}
  \end{callself}

\end{sequencediagram}
}%
\caption{Diagramme de séquence SC1 — Connexion}
\label{fig:seq-login}
\end{figure}

\subsubsection*{SC2 — Inscription (Register)}

\begin{figure}[H]
\centering
\resizebox{\textwidth}{!}{%
\begin{sequencediagram}
  \newthread{cli}{:Client}
  \newinst[1.5]{ctrl}{:AuthController}
  \newinst[1.5]{hand}{:RegisterCmd\\Handler}
  \newinst[1.5]{db}{:Database}

  \begin{call}{cli}{POST /api/auth/register\{firstName,email,password,role\}}{ctrl}{}
    \begin{call}{ctrl}{Handle(RegisterCommand)}{hand}{}
      \begin{call}{hand}{EmailExists(email)}{db}{bool}
      \end{call}
      \begin{callself}{hand}{HashPassword(password)}{}
      \end{callself}
      \begin{call}{hand}{CreateUser(user)}{db}{userId}
      \end{call}
    \end{call}
    \mess{ctrl}{200 OK \{userId, message\}}{cli}
  \end{call}

\end{sequencediagram}
}%
\caption{Diagramme de séquence SC2 — Inscription}
\label{fig:seq-register}
\end{figure}

\section{Conclusion}

Le Sprint 1 a établi l'infrastructure d'authentification complète de MecaRage.
Le système JWT est opérationnel avec gestion du refresh token, isolation multi-tenant
via filtre global EF Core, et hachage bcrypt des mots de passe. Les interfaces
Angular de connexion et d'inscription sont fonctionnelles et reliées au backend.

% ===========================================================================
% CHAPITRE 4 — SPRINT 2 : GESTION DES DONNÉES DE RÉFÉRENCE
% ===========================================================================
\chapter{Sprint 2 : Gestion des Données de Référence}

\section{Introduction}

Ce sprint met en place les données fondamentales sur lesquelles repose l'exploitation
de la plateforme~: la gestion des tenants, des garages, du personnel, des véhicules
clients et de l'inventaire des pièces détachées. Ces entités constituent le référentiel
à partir duquel toutes les opérations métier s'appuieront.

\section{Identification du Backlog Sprint 2}

\begin{table}[H]
\centering
\caption{Backlog Sprint 2 — Données de Référence}
\begin{tabularx}{\textwidth}{|l|X|l|c|}
\hline
\rowcolor{indigolite!60}
\textbf{ID} & \textbf{User Story} & \textbf{Priorité} & \textbf{SP} \\
\hline
US05 & Super Admin~: gérer les tenants (créer, modifier, supprimer) &
  \textcolor{prioH}{\bfseries Haute} & 3 \\
\hline
US06 & Admin Entreprise~: créer et gérer des garages &
  \textcolor{prioH}{\bfseries Haute} & 5 \\
\hline
US07 & Admin Entreprise~: créer et gérer le personnel &
  \textcolor{prioH}{\bfseries Haute} & 5 \\
\hline
US08 & Client~: enregistrer et gérer ses véhicules &
  \textcolor{prioH}{\bfseries Haute} & 3 \\
\hline
US09 & Admin~: gérer l'inventaire des pièces détachées &
  \textcolor{prioM}{\bfseries Moyenne} & 5 \\
\hline
\multicolumn{3}{|r|}{\textbf{Total Sprint 2}} & \textbf{21 SP} \\
\hline
\end{tabularx}
\end{table}

\noindent\textbf{Tâches techniques du Sprint 2~:}
\begin{enumerate}[noitemsep]
  \item Implémentation des entités \texttt{Garage}, \texttt{Vehicle}, \texttt{SparePart}.
  \item Commandes CQRS~: \texttt{CreateGarageCommand}, \texttt{CreateVehicleCommand},
    \texttt{CreateSparePartCommand}, \texttt{CreateStaffCommand}.
  \item Migrations EF Core pour les nouvelles tables.
  \item Développement des contrôleurs REST~: \texttt{GaragesController},
    \texttt{VehiclesController}, \texttt{SparePartsController}.
  \item Interfaces Angular~: tableau de bord admin, gestion du garage et inventaire.
  \item Application des policies d'autorisation par rôle (\texttt{[Authorize(Roles = ...)]}).
\end{enumerate}

\section{Conception}

\subsection{Diagramme de classes — Sprint 2}

\begin{figure}[H]
\centering
\resizebox{\textwidth}{!}{%
\begin{tikzpicture}[node distance=0.8cm and 3cm]

\node[umlclass] (garage) {
  \textbf{Garage}
  \nodepart{two}
  + id : Guid\\+ tenantId : Guid\\+ name : String\\+ address : String\\+ city : String\\+ phone : String\\+ latitude : double?\\+ longitude : double?\\+ adminId : Guid?\\+ isActive : bool
  \nodepart{three}
  (géré via commands)
};

\node[umlclass,right=3cm of garage] (vehicle) {
  \textbf{Vehicle}
  \nodepart{two}
  + id : Guid\\+ clientId : Guid\\+ brand : String\\+ model : String\\+ year : int\\+ licensePlate : String\\+ fuelType : String\\+ mileage : int\\+ vin : String?
  \nodepart{three}
  (géré via commands)
};

\node[umlclass,right=3cm of vehicle] (spare) {
  \textbf{SparePart}
  \nodepart{two}
  + id : Guid\\+ garageId : Guid\\+ code : String\\+ name : String\\+ category : String\\+ unitPrice : decimal\\+ stockQty : int\\+ reorderLevel : int\\+ manufacturer : String?
  \nodepart{three}
  (géré via commands)
};

\node[umlclass,below=2.8cm of garage] (user2) {
  \textbf{User (Staff)}
  \nodepart{two}
  + id : Guid\\+ firstName : String\\+ email : String\\+ role : UserRole\\+ garageId : Guid?
  \nodepart{three}
  \textit{(cf. Sprint 1)}
};

\draw[classrel] (garage.east) -- node[above,font=\scriptsize]{1}
  node[below,font=\scriptsize]{N} (vehicle.west);
\draw[classcomp] (garage.east) -- node[above,font=\scriptsize]{1}
  node[below,font=\scriptsize]{N} (spare.west);
\draw[classrel] (user2.east) -- node[above,font=\scriptsize]{N}
  node[below,font=\scriptsize]{assigné} (garage.south);

\end{tikzpicture}
}%
\caption{Diagramme de classes — Sprint 2 (Données de Référence)}
\label{fig:class-sprint2}
\end{figure}

\subsection{Diagrammes de séquence}

\subsubsection*{SC3 — Création d'un Garage}

\begin{figure}[H]
\centering
\resizebox{\textwidth}{!}{%
\begin{sequencediagram}
  \newthread{adm}{:AdminEntreprise}
  \newinst[2]{ctrl}{:GaragesController}
  \newinst[2]{hand}{:CreateGarage\\CmdHandler}
  \newinst[2]{db}{:Database}

  \begin{call}{adm}{POST /api/garages\{name,address,city,tenantId\}}{ctrl}{}
    \begin{call}{ctrl}{Handle(CreateGarageCommand)}{hand}{}
      \begin{call}{hand}{TenantExists(tenantId)}{db}{bool}
      \end{call}
      \begin{callself}{hand}{MapToEntity(command)}{}
      \end{callself}
      \begin{call}{hand}{Add(garage)}{db}{garageId}
      \end{call}
    \end{call}
    \mess{ctrl}{201 Created \{garageId, name\}}{adm}
  \end{call}

\end{sequencediagram}
}%
\caption{Diagramme de séquence SC3 — Création d'un Garage}
\label{fig:seq-garage}
\end{figure}

\subsubsection*{SC4 — Enregistrement d'un Véhicule}

\begin{figure}[H]
\centering
\resizebox{\textwidth}{!}{%
\begin{sequencediagram}
  \newthread{cli}{:Client}
  \newinst[2]{ctrl}{:VehiclesController}
  \newinst[2]{hand}{:CreateVehicle\\CmdHandler}
  \newinst[2]{db}{:Database}

  \begin{call}{cli}{POST /api/vehicles\{brand,model,year,licensePlate,fuelType\}}{ctrl}{}
    \begin{call}{ctrl}{Handle(CreateVehicleCommand)}{hand}{}
      \begin{callself}{hand}{ExtractClientId(userClaims)}{}
      \end{callself}
      \begin{call}{hand}{LicensePlateExists(plate)}{db}{bool}
      \end{call}
      \begin{callself}{hand}{MapToEntity(command)}{}
      \end{callself}
      \begin{call}{hand}{Add(vehicle)}{db}{vehicleId}
      \end{call}
    \end{call}
    \mess{ctrl}{201 Created \{vehicleId\}}{cli}
  \end{call}

\end{sequencediagram}
}%
\caption{Diagramme de séquence SC4 — Enregistrement d'un Véhicule}
\label{fig:seq-vehicle}
\end{figure}

\section{Conclusion}

Le Sprint 2 a mis en place les données de référence essentielles~: garages,
personnel, véhicules et inventaire. Les politiques d'autorisation ont été
appliquées de façon fine par rôle, assurant que seul un \texttt{AdminEntreprise}
peut créer un garage, et qu'un \texttt{Client} ne gère que ses propres véhicules.

% ===========================================================================
% CHAPITRE 5 — SPRINT 3 : RAPPORTS DE SYMPTÔMES, DIAGNOSTIC IA ET RENDEZ-VOUS
% ===========================================================================
\chapter{Sprint 3 : Rapports de Symptômes, Diagnostic IA et Rendez-vous}

\section{Introduction}

Ce sprint implémente la boucle de diagnostic préliminaire et de prise de rendez-vous.
Une fois son véhicule enregistré, le client soumet une description textuelle des
symptômes. Le service IA (FastAPI + Google Gemini) génère automatiquement un diagnostic
structuré avec un score de confiance. Le chef d'atelier révise ensuite ce rapport
et propose une plage de disponibilité pour le rendez-vous.

\section{Identification du Backlog Sprint 3}

\begin{table}[H]
\centering
\caption{Backlog Sprint 3 — Symptômes, IA et Rendez-vous}
\begin{tabularx}{\textwidth}{|l|X|l|c|}
\hline
\rowcolor{indigolite!60}
\textbf{ID} & \textbf{User Story} & \textbf{Priorité} & \textbf{SP} \\
\hline
US10 & Client~: soumettre un rapport de symptômes textuels &
  \textcolor{prioH}{\bfseries Haute} & 5 \\
\hline
US11 & Système~: générer un diagnostic IA (RAG + Gemini) automatiquement &
  \textcolor{prioH}{\bfseries Haute} & 8 \\
\hline
US12 & Chef d'Atelier~: réviser les rapports et proposer une plage RDV &
  \textcolor{prioH}{\bfseries Haute} & 5 \\
\hline
US13 & Client~: réserver un rendez-vous dans la plage proposée &
  \textcolor{prioH}{\bfseries Haute} & 3 \\
\hline
US14 & Chef d'Atelier~: approuver ou refuser le rendez-vous en assignant un mécanicien &
  \textcolor{prioH}{\bfseries Haute} & 3 \\
\hline
\multicolumn{3}{|r|}{\textbf{Total Sprint 3}} & \textbf{24 SP} \\
\hline
\end{tabularx}
\end{table}

\section{Conception}

\subsection{Diagramme de classes — Sprint 3}

\begin{figure}[H]
\centering
\resizebox{\textwidth}{!}{%
\begin{tikzpicture}[node distance=0.8cm and 3.5cm]

\node[umlclass] (sym) {
  \textbf{SymptomReport}
  \nodepart{two}
  + id : Guid\\+ clientId : Guid\\+ vehicleId : Guid\\+ garageId : Guid?\\+ symptomsDescription : String\\+ aiPredictedIssue : String?\\+ aiConfidenceScore : float?\\+ aiRecommendations : String?\\+ chefFeedback : String?\\+ status : SymptomReportStatus\\+ availablePeriodStart : DateTime?\\+ availablePeriodEnd : DateTime?
  \nodepart{three}
  (géré via commands)
};

\node[umlclass,right=4cm of sym] (appt) {
  \textbf{Appointment}
  \nodepart{two}
  + id : Guid\\+ clientId : Guid\\+ vehicleId : Guid\\+ garageId : Guid\\+ symptomReportId : Guid?\\+ preferredDate : DateTime\\+ preferredTime : TimeSpan\\+ specialRequests : String?\\+ status : AppointmentStatus\\+ approvedByChefId : Guid?\\+ approvedAt : DateTime?
  \nodepart{three}
  (géré via commands)
};

\node[umlclass,below=2.5cm of sym] (aisvc) {
  \textbf{IADiagnosisService}
  \nodepart{two}
  + endpoint : String\\+ model : String (Gemini)\\+ confidence : float
  \nodepart{three}
  + diagnose(symptoms) : Result
};

\draw[classrel] (appt.west) -- node[above,font=\scriptsize]{1}
  node[below,font=\scriptsize]{0..1} (sym.east)
  node[pos=0.5,below=2pt,font=\scriptsize]{liée à};

\draw[classrel,dashed] (sym.south) -- node[right,font=\scriptsize]
  {\texttt{<<use>>}} (aisvc.north);

\end{tikzpicture}
}%
\caption{Diagramme de classes — Sprint 3 (Symptômes, IA, Rendez-vous)}
\label{fig:class-sprint3}
\end{figure}

\subsection{Pipeline du Diagnostic IA}

Le service IA utilise une architecture RAG (Retrieval Augmented Generation)~:

\begin{enumerate}[noitemsep]
  \item \textbf{BM25 Retrieval}~: les symptômes soumis par le client sont comparés
    à une base de connaissances mécanique (\texttt{knowledge\_base.txt}) via l'algorithme BM25.
  \item \textbf{Contexte enrichi}~: les K passages les plus pertinents sont injectés
    dans le prompt envoyé à Google Gemini.
  \item \textbf{Génération structurée}~: Gemini retourne un JSON structuré contenant
    le problème probable (\texttt{aiPredictedIssue}), un score de confiance
    (\texttt{aiConfidenceScore}) et des recommandations d'action.
\end{enumerate}

\subsection{Diagrammes de séquence}

\subsubsection*{SC5 — Soumission d'un Rapport de Symptômes avec Diagnostic IA}

\begin{figure}[H]
\centering
\resizebox{\textwidth}{!}{%
\begin{sequencediagram}
  \newthread{cli}{:Client}
  \newinst[1.5]{sctrl}{:SymptomReports\\Controller}
  \newinst[1.5]{hand}{:CreateSymptom\\ReportCmdHandler}
  \newinst[1.5]{ia}{:IAService\\(FastAPI)}
  \newinst[1.5]{db}{:Database}
  \newinst[1.5]{notif}{:Notification\\Service}

  \begin{call}{cli}{POST /api/symptom-reports\{vehicleId,symptoms\}}{sctrl}{}
    \begin{call}{sctrl}{Handle(CreateSymptomReportCommand)}{hand}{}
      \begin{call}{hand}{POST /diagnose\{symptoms\}}{ia}{DiagnosisResult}
      \end{call}
      \begin{callself}{hand}{MapDiagnosisToReport()}{}
      \end{callself}
      \begin{call}{hand}{Add(symptomReport)}{db}{reportId}
      \end{call}
      \begin{call}{hand}{NotifyChefAtelier(reportId)}{notif}{}
        \begin{call}{notif}{Add(notification)}{db}{}
        \end{call}
      \end{call}
    \end{call}
    \mess{sctrl}{201 Created \{reportId, aiDiagnosis\}}{cli}
  \end{call}

\end{sequencediagram}
}%
\caption{Diagramme de séquence SC5 — Soumission d'un Rapport de Symptômes}
\label{fig:seq-symptom}
\end{figure}

\subsubsection*{SC6 — Approbation d'un Rendez-vous par le Chef d'Atelier}

\begin{figure}[H]
\centering
\resizebox{\textwidth}{!}{%
\begin{sequencediagram}
  \newthread{chef}{:ChefAtelier}
  \newinst[2]{ctrl}{:Appointments\\Controller}
  \newinst[2]{hand}{:ApproveAppt\\CmdHandler}
  \newinst[2]{db}{:Database}
  \newinst[2]{notif}{:Notification\\Service}
  \newinst[2]{ihand}{:CreateIntervention\\LifecycleCmdHandler}

  \begin{call}{chef}{PATCH /api/appointments/\{id\}/approve}{ctrl}{}
    \begin{call}{ctrl}{Handle(ApproveAppointmentCommand)}{hand}{}
      \begin{call}{hand}{GetAppointment(id)}{db}{Appointment}
      \end{call}
      \begin{callself}{hand}{SetStatus(Approved)}{}
      \end{callself}
      \begin{call}{hand}{SaveChanges()}{db}{}
      \end{call}
      \begin{call}{hand}{CreateInterventionLifecycle()}{ihand}{}
        \begin{call}{ihand}{Add(Intervention)}{db}{}
        \end{call}
      \end{call}
      \begin{call}{hand}{NotifyClient(appointmentId)}{notif}{}
        \begin{call}{notif}{Add(notification)}{db}{}
        \end{call}
      \end{call}
    \end{call}
    \mess{ctrl}{200 OK}{chef}
  \end{call}

\end{sequencediagram}
}%
\caption{Diagramme de séquence SC6 — Approbation d'un Rendez-vous}
\label{fig:seq-approve-appt}
\end{figure}

\section{Conclusion}

Le Sprint 3 a réalisé la première boucle de valeur visible par le client~: soumettre
des symptômes, recevoir un diagnostic IA instantané, obtenir une plage de rendez-vous,
et être notifié de sa confirmation. L'approbation d'un rendez-vous déclenche
automatiquement la création du traceur de cycle de vie de l'intervention.

% ===========================================================================
% CHAPITRE 6 — SPRINT 4 : OPÉRATIONS D'ATELIER ET CYCLE DE VIE
% ===========================================================================
\chapter{Sprint 4 : Opérations d'Atelier et Cycle de Vie des Interventions}

\section{Introduction}

Ce sprint représente le c\oe ur des opérations de l'atelier. Il couvre l'affectation
des mécaniciens, la remontée des rapports d'examen, la facturation et son approbation
par le client, l'exécution et la validation de la réparation, la génération PDF de la
facture, et enfin l'enregistrement du paiement. Toutes ces étapes sont automatiquement
tracées par le \textbf{Intervention Lifecycle Tracker}, une entité transverse qui
maintient l'état global d'une intervention à travers 10 états distincts.

\section{Identification du Backlog Sprint 4}

\begin{table}[H]
\centering
\caption{Backlog Sprint 4 — Opérations d'Atelier}
\begin{tabularx}{\textwidth}{|l|X|l|c|}
\hline
\rowcolor{indigolite!60}
\textbf{ID} & \textbf{User Story} & \textbf{Priorité} & \textbf{SP} \\
\hline
US15 & Chef~: assigner mécaniciens à une tâche de réparation &
  \textcolor{prioH}{\bfseries Haute} & 5 \\
\hline
US16 & Mécanicien~: consulter ses tâches assignées &
  \textcolor{prioH}{\bfseries Haute} & 3 \\
\hline
US17 & Mécanicien~: soumettre rapport d'examen (pièces, observations) &
  \textcolor{prioH}{\bfseries Haute} & 5 \\
\hline
US18 & Chef~: réviser et valider le rapport d'examen &
  \textcolor{prioH}{\bfseries Haute} & 3 \\
\hline
US19 & Chef~: créer et finaliser un devis (facture) &
  \textcolor{prioH}{\bfseries Haute} & 5 \\
\hline
US20 & Client~: approuver ou refuser le devis &
  \textcolor{prioH}{\bfseries Haute} & 3 \\
\hline
US21 & Mécanicien~: soumettre la fin de réparation &
  \textcolor{prioH}{\bfseries Haute} & 3 \\
\hline
US22 & Chef~: valider réparation et notifier le client (véhicule prêt) &
  \textcolor{prioH}{\bfseries Haute} & 3 \\
\hline
US23 & Admin~: enregistrer le paiement (cash/carte/virement) &
  \textcolor{prioH}{\bfseries Haute} & 3 \\
\hline
US24 & Système~: tracer le cycle de vie automatiquement &
  \textcolor{prioH}{\bfseries Haute} & 8 \\
\hline
US25 & Utilisateur~: recevoir notifications in-app et push (FCM) &
  \textcolor{prioM}{\bfseries Moyenne} & 5 \\
\hline
US26 & Client/Mécanicien~: télécharger la facture en PDF &
  \textcolor{prioM}{\bfseries Moyenne} & 3 \\
\hline
\multicolumn{3}{|r|}{\textbf{Total Sprint 4}} & \textbf{49 SP} \\
\hline
\end{tabularx}
\end{table}

\subsection*{Machine d'états — Intervention Lifecycle}

La figure~\ref{fig:lifecycle-states} présente les 10 transitions du suivi d'intervention~:

\begin{figure}[H]
\centering
\begin{tikzpicture}[
  state/.style={rectangle,draw=indigo,fill=indigolite!30,rounded corners=4pt,
    font=\scriptsize\bfseries,text width=2.8cm,align=center,minimum height=0.7cm},
  trans/.style={-{Stealth},draw=indigo!70,font=\scriptsize},
  node distance=1.2cm]

  \node[state] (c)   {Created};
  \node[state,below=of c] (ue)  {UnderExamination};
  \node[state,below=of ue] (er)  {ExaminationReviewed};
  \node[state,below=of er] (ip)  {InvoicePending};
  \node[state,below left=1cm and 1.5cm of ip] (ap)  {Approved\\(RepairInProgress)};
  \node[state,below right=1cm and 1.5cm of ip] (rej) {Rejected};
  \node[state,below=of ap] (rc)  {RepairCompleted};
  \node[state,below=of rc] (rfp) {ReadyForPickup};
  \node[state,below=of rfp] (cl)  {Closed};

  \draw[trans] (c)   -- node[right,font=\scriptsize]{ExamSubmit} (ue);
  \draw[trans] (ue)  -- node[right,font=\scriptsize]{ChefReview} (er);
  \draw[trans] (er)  -- node[right,font=\scriptsize]{InvoiceFinalized} (ip);
  \draw[trans] (ip)  -- node[left,font=\scriptsize]{ClientApproves} (ap);
  \draw[trans] (ip)  -- node[right,font=\scriptsize]{ClientRejects} (rej);
  \draw[trans] (ap)  -- node[right,font=\scriptsize]{RepairSubmit} (rc);
  \draw[trans] (rc)  -- node[right,font=\scriptsize]{ChefValidates} (rfp);
  \draw[trans] (rfp) -- node[right,font=\scriptsize]{AdminPayment} (cl);

\end{tikzpicture}
\caption{Machine d'états — Intervention Lifecycle (10 états)}
\label{fig:lifecycle-states}
\end{figure}

\section{Conception}

\subsection{Diagramme de classes — Sprint 4}

\begin{figure}[H]
\centering
\resizebox{\textwidth}{!}{%
\begin{tikzpicture}[node distance=0.8cm and 3cm]

\node[umlclass] (rt) {
  \textbf{RepairTask}
  \nodepart{two}
  + id : Guid\\+ appointmentId : Guid\\+ taskTitle : String\\+ description : String\\+ status : RepairTaskStatus\\+ assignedAt : DateTime\\+ completionNotes : String?\\+ estimatedMinutes : int?
  \nodepart{three}
  (géré via commands)
};

\node[umlclass,right=2.5cm of rt] (rta) {
  \textbf{RepairTaskAssignment}
  \nodepart{two}
  + id : Guid\\+ repairTaskId : Guid\\+ mechanicId : Guid\\+ examinationObs : String?\\+ partsNeeded : String?\\+ examinationStatus : String
  \nodepart{three}
  (géré via commands)
};

\node[umlclass,right=2.5cm of rta] (inv) {
  \textbf{Invoice}
  \nodepart{two}
  + id : Guid\\+ appointmentId : Guid\\+ invoiceNumber : String\\+ serviceFee : decimal\\+ partsTotal : decimal\\+ totalAmount : decimal\\+ status : InvoiceStatus\\+ clientApproved : bool
  \nodepart{three}
  (géré via commands)
};

\node[umlclass,below=2.5cm of rta] (interv) {
  \textbf{Intervention}\\(Lifecycle Tracker)
  \nodepart{two}
  + id : Guid\\+ appointmentId : Guid?\\+ invoiceId : Guid?\\+ repairTaskId : Guid?\\+ status : InterventionLifecycleStatus\\+ proceedWithIntervention : bool?\\+ examinationNotes : String?\\+ repairNotes : String?\\+ paymentAmount : decimal?\\+ paymentMethod : String?
  \nodepart{three}
  (side tracker automatique)
};

\node[umlclass,left=2.5cm of interv] (ili) {
  \textbf{InvoiceLineItem}
  \nodepart{two}
  + id : Guid\\+ invoiceId : Guid\\+ sparePartId : Guid?\\+ description : String\\+ quantity : int\\+ unitPrice : decimal\\+ lineTotal : decimal
  \nodepart{three}
  (géré via commands)
};

\draw[classcomp] (rt.east) -- node[above,font=\scriptsize]{1}
  node[below,font=\scriptsize]{N} (rta.west);
\draw[classcomp] (inv.south) -- node[right,font=\scriptsize]{1 N}
  (ili.north east);
\draw[classrel] (rt.south) -- node[right,font=\scriptsize]{}(interv.north west);
\draw[classrel] (inv.south) -- node[left,font=\scriptsize]{}(interv.north east);

\end{tikzpicture}
}%
\caption{Diagramme de classes — Sprint 4 (Atelier et Cycle de Vie)}
\label{fig:class-sprint4}
\end{figure}

\subsection{Diagrammes de séquence}

\subsubsection*{SC7 — Finalisation de la Facture et Approbation Client}

\begin{figure}[H]
\centering
\resizebox{\textwidth}{!}{%
\begin{sequencediagram}
  \newthread{chef}{:ChefAtelier}
  \newinst[1.5]{ictrl}{:Invoices\\Controller}
  \newinst[1.5]{fhand}{:FinalizeInvoice\\CmdHandler}
  \newinst[1.5]{db}{:Database}
  \newinst[1.5]{notif}{:Notification\\Service}
  \newthread{cli}{:Client}

  \begin{call}{chef}{PATCH /api/invoices/\{id\}/finalize}{ictrl}{}
    \begin{call}{ictrl}{Handle(FinalizeInvoiceCommand)}{fhand}{}
      \begin{call}{fhand}{GetInvoice(id)}{db}{Invoice}
      \end{call}
      \begin{callself}{fhand}{SetStatus(AwaitingApproval)}{}
      \end{callself}
      \begin{call}{fhand}{UpdateInterventionStatus(InvoicePending)}{db}{}
      \end{call}
      \begin{call}{fhand}{NotifyClient(invoiceId)}{notif}{}
        \begin{call}{notif}{Add(notification)}{db}{}
        \end{call}
      \end{call}
    \end{call}
    \mess{ictrl}{200 OK}{chef}
  \end{call}

  \mess{notif}{Notification: Devis disponible}{cli}

  \begin{call}{cli}{POST /api/invoices/\{id\}/approve}{ictrl}{}
    \begin{call}{ictrl}{Handle(ApproveInvoiceCommand)}{fhand}{}
      \begin{callself}{fhand}{SetStatus(Approved), ProceedWithIntervention=true}{}
      \end{callself}
      \begin{call}{fhand}{UpdateInterventionStatus(RepairInProgress)}{db}{}
      \end{call}
    \end{call}
    \mess{ictrl}{200 OK}{cli}
  \end{call}

\end{sequencediagram}
}%
\caption{Diagramme de séquence SC7 — Finalisation et Approbation du Devis}
\label{fig:seq-invoice}
\end{figure}

\subsubsection*{SC8 — Enregistrement du Paiement et Clôture de l'Intervention}

\begin{figure}[H]
\centering
\resizebox{\textwidth}{!}{%
\begin{sequencediagram}
  \newthread{adm}{:AdminEntreprise}
  \newinst[2]{ictrl}{:Interventions\\Controller}
  \newinst[2]{phand}{:RegisterPayment\\CmdHandler}
  \newinst[2]{db}{:Database}
  \newinst[2]{notif}{:Notification\\Service}

  \begin{call}{adm}{PATCH /api/interventions/lifecycle/\{id\}/payment}{ictrl}{}
    \begin{call}{ictrl}{Handle(RegisterPaymentCommand)}{phand}{}
      \begin{call}{phand}{GetIntervention(id)}{db}{Intervention}
      \end{call}
      \begin{callself}{phand}{VerifyStatus(ReadyForPickup)}{}
      \end{callself}
      \begin{callself}{phand}{SetPaymentFields + Status=Closed}{}
      \end{callself}
      \begin{call}{phand}{UpdateInvoiceStatus(Paid)}{db}{}
      \end{call}
      \begin{call}{phand}{SaveChanges()}{db}{}
      \end{call}
    \end{call}
    \mess{ictrl}{200 OK \{status: Closed\}}{adm}
  \end{call}

\end{sequencediagram}
}%
\caption{Diagramme de séquence SC8 — Enregistrement du Paiement}
\label{fig:seq-payment}
\end{figure}

\section{Conclusion}

Le Sprint 4 a réalisé l'ensemble des opérations métier d'un atelier mécanique~:
de l'affectation des mécaniciens jusqu'à la clôture de l'intervention avec paiement.
Le Intervention Lifecycle Tracker automatiquement mis à jour à chaque étape offre
une traçabilité complète sur les 10 états du cycle de vie. La génération de facture
PDF (QuestPDF) et les notifications FCM (Firebase) complètent l'expérience utilisateur
sur toutes les plateformes (web Angular et mobile Android).

% ===========================================================================
% CONCLUSION GÉNÉRALE
% ===========================================================================
\chapter*{Conclusion Générale}
\addcontentsline{toc}{chapter}{Conclusion Générale}

\setlength{\parskip}{6pt}

Le projet \textbf{MecaRage} représente une solution logicielle complète et cohérente
pour la modernisation de la gestion des ateliers mécaniques. À travers quatre sprints
successifs, nous avons construit une plateforme qui couvre l'intégralité du cycle de
vie d'une intervention automobile~: de la réception initiale des symptômes jusqu'au
paiement final.

\bigskip

Les principaux apports techniques du projet sont~:
\begin{itemize}[noitemsep]
  \item Une \textbf{architecture Clean Architecture avec CQRS} garantissant la
    séparation des préoccupations, la testabilité et la maintenabilité à long terme.
  \item Un \textbf{service de diagnostic IA} basé sur une approche RAG (BM25 + Google
    Gemini) permettant un pré-diagnostic automatisé à partir d'une simple description
    textuelle des symptômes.
  \item Un \textbf{Intervention Lifecycle Tracker} — suivi transverse automatique
    traversant 10 états, déclenché par les événements existants sans modifier le flux
    opérationnel.
  \item Une \textbf{application mobile Android} (Java) connectée via Firebase pour
    les notifications temps réel (RTDB + FCM), avec la plateforme REST comme
    source de vérité unique.
  \item Un \textbf{système de génération PDF} (QuestPDF) pour les factures
    professionnelles téléchargeables directement depuis l'interface client.
  \item Une \textbf{infrastructure de déploiement} conteneurisée via Docker Compose,
    supervisée par Prometheus et Grafana.
\end{itemize}

\bigskip

Les perspectives d'évolution de MecaRage incluent~: l'enrichissement de la base de
connaissances du service de diagnostic IA, l'implémentation d'alertes de réapprovisionnement
automatique des pièces détachées, l'intégration d'un module de gestion financière avec
export comptable, et l'ajout d'un calendrier interactif pour la planification des
rendez-vous.

\bigskip

Ce projet nous a permis d'appliquer concrètement des patterns d'architecture logicielle
modernes, d'intégrer des technologies hétérogènes (.NET, Angular, Python, Java, Firebase)
dans une solution cohérente, et de répondre à des besoins métier réels dans un
contexte multi-tenant contraignant.

% ===========================================================================
% BIBLIOGRAPHIE
% ===========================================================================
\chapter*{Références}
\addcontentsline{toc}{chapter}{Références}

\begin{enumerate}[label={[\arabic*]}]
  \item Microsoft. \textit{ASP.NET Core 9 Documentation}. Microsoft Docs, 2025.
    \url{https://docs.microsoft.com/aspnet/core}
  \item Google. \textit{Firebase Realtime Database Documentation}. Google Developers, 2025.
    \url{https://firebase.google.com/docs/database}
  \item Google. \textit{Gemini API Documentation}. Google AI Studio, 2025.
    \url{https://ai.google.dev}
  \item Angular Team. \textit{Angular v21 Documentation}. Angular.io, 2025.
    \url{https://angular.io/docs}
  \item Fowler, M. \textit{Patterns of Enterprise Application Architecture}. Addison-Wesley, 2002.
  \item Young, G. \textit{CQRS Documents}. 2010.
    \url{https://cqrs.files.wordpress.com/2010/11/cqrs_documents.pdf}
  \item QuestPDF. \textit{QuestPDF Documentation}. 2025.
    \url{https://www.questpdf.com}
  \item Robertson, S. \& Robertson, J. \textit{Mastering the Requirements Process}. Addison-Wesley, 2012.
\end{enumerate}

\end{document}

